%% Version 7.0.1. Created May 2025.
%% Version 7. Created January 2025.

\documentclass[twocolumn,twocolappendix,linenumbers]{aastex701}
\usepackage{amsmath}
\usepackage{amssymb}
\usepackage{booktabs}
\usepackage{xcolor}
\newcommand{\N}{\mathcal{N}}
\newcommand{\ind}{\mathbf{1}}
\newcommand{\dd}{\mathrm{d}}
\newcommand{\cat}{\mathrm{cat}}
\newcommand{\bg}{\mathrm{bg}}
\newcommand{\lens}{\mathrm{lens}}
\newcommand{\logit}{\operatorname{logit}}

\newcommand{\isaque}[1]{\textcolor{teal}{[Isaque: #1]}}
\newcommand{\priya}[1]{\textcolor{orange}{[Priya: #1]}}
\newcommand{\colin}[1]{\textcolor{magenta}{[Colin: #1]}}

\newcommand{\todo}[1]{\textcolor{red!65!black}{[To complete: #1]}}
\newcommand{\rev}[1]{\textcolor{red}{#1}}
\newcommand{\lenstool}{\textsc{Lenstool}\xspace}

\graphicspath{{./}{figures/}}
\begin{document}
\title{Probabilistic Assignment of Multiple Images in Galaxy Cluster Lenses}

\author[0000-0001-9947-6911]{Colin J. Burke}
\affiliation{Department of Physics, University of North Texas, Denton, TX 76201, USA}
\affiliation{Department of Astronomy, Yale University, New Haven, CT 06511, USA}
\email{colin.j.burke@yale.edu}

\author[0000-0001-7040-4930]{Isaque Dutra}
\affiliation{Department of Physics, Yale University, New Haven, CT 06511, USA}
\email{isaque.dutra@yale.edu}
\correspondingauthor{Isaque Dutra}

\author[0000-0002-5554-8896]{Priyamvada Natarajan}
\affiliation{Department of Physics, Yale University, New Haven, CT 06511, USA}
\affiliation{Department of Astronomy, Yale University, New Haven, CT 06511, USA}
\email{priyamvada.natarajan@yale.edu}

\begin{abstract}
\end{abstract}

\keywords{\uat{Strong gravitational lensing}{1643} --- \uat{Galaxy clusters}{584} --- \uat{Dark matter distribution}{356} --- \uat{Galaxy dark matter halos}{1880} --- \uat{Bayesian statistics}{1900} --- \uat{Astronomy software}{1855}}

\begin{abstract}
Identifying which images share a background source is part of constructing a cluster strong-lensing model.
When more than one association is plausible, fixing the image families removes uncertainty that can affect the inferred mass distribution.
We describe a method that combines catalog measurements and source-plane consistency to weight competing assignments.
Each assignment is a complete partition of the observed images, so an image belongs to exactly one family in each candidate.
Catalog scores compare a shared latent property with independent background draws, while the lensing score measures agreement around a fitted source position.
We normalize the combined scores over a generated set of partitions and obtain pairwise same-family probabilities by summing their weights.
The resulting distribution preserves consistent assignments but remains conditional on the candidate search, adopted scores, and lens model.
This draft specifies the implemented method and the validation needed to assess recovery and probability calibration.
\end{abstract}

\section{Introduction}
\label{sec:introduction}
Cluster strong-lensing models use the positions and redshifts of multiple images to constrain the projected mass distribution \citep{Jullo2007}.
Images assigned to one family must be consistent with a common background source.
An incorrect association can therefore force the mass model to explain an image configuration that the lens did not produce.
Similar colors or uncertain redshifts can leave more than one association plausible.
We ask how catalog information and lensing consistency can retain this ambiguity through probabilities for alternative family assignments.

Comparing a common-source hypothesis with independent-source hypotheses is established in astronomical cross-identification \citep{Budavari2008}.
Requiring mutually compatible associations also motivates inference over complete catalog assignments \citep{Budavari2016}.
Here the images can be widely separated on the sky, and their positional consistency is evaluated after applying a lens model.
We describe the family-level scoring configuration of an implemented partition engine.
Its output is a weighted set of assignments and the probability, under those weights, that any two images share a source.

\section{Probabilistic Family Assignment}
\label{sec:method}
\subsection{Observations and candidate families}
\label{sec:data}
An ambiguous image must be allowed to change families without belonging to two families simultaneously.
We therefore assign the $N$ observed images to a partition $Z=\{F_1,\ldots,F_K\}$, where $K$ is the number of nonempty families and $F_k$ is a set of image indices.
The sets are disjoint and cover all images; their sizes $n_k=|F_k|$ satisfy $\sum_k n_k=N$.
A singleton is an image with no assigned counterpart in the catalog, not necessarily a physically singly imaged source.

For image $i$, the inputs are its two-dimensional angular position $\boldsymbol{x}_i$, adopted positional scale $\sigma_i$, redshift estimate $\hat z_i$ with uncertainty $s_{z,i}$, and available catalog measurements.
Image and source-plane coordinates and their uncertainty scales are in arcseconds.
The vector $\boldsymbol{\eta}$ specifies a supplied lens model and is held fixed during the primary assignment calculation.
Each family hypothesis is scored using the available measurements, with missing values omitted.

\subsection{Catalog evidence for a shared source}
\label{sec:catalog}
Images of the same source should have consistent redshifts and, under an achromatic magnification approximation, similar spectral shapes.
For a family $F$ and scalar catalog component $a$, let $y_{ia}$ be its value, $v_{ia}$ its adopted variance, and $\mu_{Fa}$ the unknown common value.
Conditional independence gives the joint density
\begin{equation}
 \mathcal L_{Fa}(\mu_{Fa})=
 \prod_{i\in F_a}\N(y_{ia}\mid\mu_{Fa},v_{ia}),
 \label{eq:catalogjoint}
\end{equation}
where $F_a\subseteq F$ contains the valid measurements and $\N(y\mid\mu,v)$ is a normalized scalar Gaussian with mean $\mu$ and variance $v$.
We use a proper prior $\mu_{Fa}\sim\N(m_a,\tau_a^2)$, where $m_a$ and $\tau_a$ specify the population location and prior width.
Integrating the shared value gives
\begin{equation}
 E_{Fa}=\int\dd\mu_{Fa}\,
 \N(\mu_{Fa}\mid m_a,\tau_a^2)
 \prod_{i\in F_a}\N(y_{ia}\mid\mu_{Fa},v_{ia}).
 \label{eq:catalogmarginal}
\end{equation}
Appendix~\ref{app:gaussian} evaluates this integral with its normalization intact.

To determine whether agreement exceeds chance similarity, we compare Equation~\ref{eq:catalogmarginal} with independent draws from a background density,
\begin{equation}
 B_{Fa}=\prod_{i\in F_a}\N(y_{ia}\mid b_a,t_a^2),
 \qquad R_{Fa}=E_{Fa}/B_{Fa},
 \label{eq:catalogratio}
\end{equation}
where $b_a$ and $t_a$ are the background mean and width.
The background width describes the observed catalog population; measurement variances are not added to it in this implementation.
For fixed hyperparameters, $R_{Fa}$ is a Bayes factor between these scalar models.
The implementation sets $\log R_{Fa}=0$ when fewer than two valid values remain, including singletons.
This is a scoring convention, not the general one-observation limit of Equation~\ref{eq:catalogratio}.

The catalog score combines redshift, centered magnitudes, and available morphology components,
\begin{equation}
 C_F=w_z\log R_{Fz}
       +w_c\sum_b\log R_{Fb}
       +w_m\sum_h\log R_{Fh}.
 \label{eq:catalogscore}
\end{equation}
Here $b$ indexes photometric bands, $h$ indexes size and ellipticity, and $w_z,w_c,w_m$ are fixed dimensionless score weights.
The centered magnitudes remove an overall brightness offset (Appendix~\ref{app:catalogdetails}).
They are correlated, although their scalar scores are added separately.
Observed size and ellipticity also change under lensing, so their agreement is a heuristic contribution rather than a shared intrinsic-property constraint.
Together with empirical background estimates and adjustable powers, these choices make $C_F$ a composite score.

\subsection{Lensing consistency}
\label{sec:lensing}
A lens model supplies an additional test because images of one source should back-project to a common position.
For trial source redshift $z$, the back-projected position is
\begin{equation}
 \widehat{\boldsymbol{\beta}}_i(\boldsymbol{\eta},z)
 =\boldsymbol{x}_i-
 \boldsymbol{\alpha}(\boldsymbol{x}_i;\boldsymbol{\eta},z),
 \label{eq:lens}
\end{equation}
where $\boldsymbol{\alpha}$ is the reduced angular deflection.
Source-plane consistency has also been used in cluster lens optimization \citep{Jullo2007}.
Here the implemented variance is $u_i=\sigma_i^2+\sigma_{\rm src}^2$, where the fixed scale $\sigma_{\rm src}$ allows additional source-plane scatter.
The positional input $\sigma_i$ is used directly in this variance.
This isotropic prescription does not propagate image-plane errors through the lens Jacobian.

With $\boldsymbol{\beta}_F$ the shared source position and $\boldsymbol I_2$ the two-dimensional identity matrix, the conditional density is
\begin{equation}
 \mathcal L_F(\boldsymbol{\beta}_F;z)
 =\prod_{i\in F}\N_2\!\left(
 \widehat{\boldsymbol{\beta}}_i\mid
 \boldsymbol{\beta}_F,u_i\boldsymbol{I}_2\right),
 \label{eq:lensjoint}
\end{equation}
where $\N_2$ denotes a normalized bivariate Gaussian.
The code maximizes this density over the source position instead of marginalizing it.
The maximizing centroid and evaluated log density are
\begin{equation}
 \begin{aligned}
 \widehat{\boldsymbol{\beta}}_F
 &=\frac{\sum_{i\in F}u_i^{-1}\widehat{\boldsymbol{\beta}}_i}
          {\sum_{i\in F}u_i^{-1}},\\
 \ell_F^{\rm prof}(z)
 &=-\sum_{i\in F}\left[
 \frac{\|\widehat{\boldsymbol{\beta}}_i-
             \widehat{\boldsymbol{\beta}}_F\|^2}{2u_i}
 +\log(2\pi u_i)\right].
 \end{aligned}
 \label{eq:profile}
\end{equation}
For equal variances the centroid is the arithmetic mean.
The density normalization is retained, but profiling does not supply the prior-volume factor required for a marginalized evidence.
When image reliabilities are used, the mixture in Appendix~\ref{app:optional} replaces this Gaussian log density while retaining the same centroid.

Chance alignments in the source plane are represented by an empirical Gaussian density $p_{\bg}(\boldsymbol{\beta}\mid z)$ fitted to all back-projected catalog positions at $z$.
The log likelihood-ratio score is
\begin{equation}
 L_F(z)=\ell_F^{\rm prof}(z)
       -\sum_{i\in F}\log p_{\bg}(\widehat{\boldsymbol{\beta}}_i\mid z).
 \label{eq:lensratio}
\end{equation}
The background uses separate coordinate variances, as specified in Appendix~\ref{app:catalogdetails}.
Equation~\ref{eq:lensratio} is dimensionless but is not a marginalized Bayes factor.

An uncertain redshift changes the deflection, so the cached calculation evaluates a fixed, equally spaced grid $\{z_g\}_{g=1}^{G}$.
Let $F_z$ contain the family members with valid redshift estimates.
The normalized grid weights and averaged lens score are
\begin{equation}
 \begin{aligned}
 \pi_{Fg}&=
 \frac{\exp[-\tfrac12\sum_{i\in F_z}
           (z_g-\hat z_i)^2/s_{z,i}^2]}
      {\sum_{r=1}^{G}\exp[-\tfrac12\sum_{i\in F_z}
           (z_r-\hat z_i)^2/s_{z,i}^2]},\\
 L_F&=\log\sum_{g=1}^{G}\pi_{Fg}\exp[L_F(z_g)].
 \end{aligned}
 \label{eq:redshift}
\end{equation}
With no redshift information, $\pi_{Fg}=1/G$.
This is a weighted average of ratios, not generally a ratio of marginalized evidences.
The same redshift measurements also enter $C_F$; this use is part of the composite construction, not a demonstrated factorization of a joint generative model.
Single-image families contribute $L_F=0$.

\subsection{Structural preferences and assignment weights}
\label{sec:weights}
Weak evidence can otherwise favor implausible fragmentation or very large families.
The code therefore uses a partition weight with a Chinese-restaurant-process factor \citep{Pitman1995} and additional size penalties,
\begin{equation}
 \begin{aligned}
 \log\widetilde p_{\rm str}(Z)
 ={}&K\log a_0+\log\Gamma(a_0)-\log\Gamma(a_0+N)\\
 &+\sum_{k=1}^{K}\log\Gamma(n_k)
 -\lambda_1\sum_{k=1}^{K}\ind(n_k=1)\\
 &-\frac12\sum_{k:n_k\geq2}
       \left(\frac{n_k-n_0}{s_n}\right)^2.
 \end{aligned}
 \label{eq:prior}
\end{equation}
Here $\Gamma$ is the gamma function, $a_0>0$ controls the base partition distribution, $\lambda_1$ penalizes singletons, and $n_0,s_n$ set the preferred multi-image size and its tolerance.
The indicator $\ind(\cdot)$ equals one when its argument holds and zero otherwise.
The tilde denotes an unnormalized weight after adding the penalties; its normalization cancels for fixed hyperparameters.

Trusted associations and incompatible pairs enter through anchors.
For anchored pair $(i,j)\in\mathcal A$, let $a_{ij}\in\{0,1\}$ denote the required different- or same-family relation and $r_{ij}>0$ its strength.
Writing $S_{ij}(Z)=\ind(i,j\text{ belong to the same }F_k)$, the penalty is
\begin{equation}
 A(Z)=-\kappa\sum_{(i,j)\in\mathcal A}
 r_{ij}\ind[S_{ij}(Z)\ne a_{ij}],
 \label{eq:anchors}
\end{equation}
with fixed $\kappa=20$.
This is a finite penalty, not a hard constraint.
Anchors derived from the same colors or redshifts used in the scores are additional data-dependent preferences.

We combine the family scores and structural terms as
\begin{equation}
 s(Z)=\log\widetilde p_{\rm str}(Z)+A(Z)
       +\sum_{F\in Z}\left[\omega_c C_F+\omega_\ell L_F\right],
 \label{eq:score}
\end{equation}
where $\omega_c$ and $\omega_\ell$ are fixed dimensionless weights for catalog and lensing information.
No independent pair-likelihood term is included in this configuration.
For the distinct generated candidates $\mathcal C=\{Z_1,\ldots,Z_M\}$, their weights are
\begin{equation}
 W_m=\frac{\exp[s(Z_m)/T]}{\sum_{q=1}^{M}\exp[s(Z_q)/T]},
 \qquad T>0,
 \label{eq:weights}
\end{equation}
where $M$ is the candidate count and $T$ is a fixed dimensionless temperature.
Increasing $T$ flattens the distribution and tempers the structural terms as well as the data scores.
The weights sum to one over $\mathcal C$; they do not measure omitted probability outside it.

To retain uncertainty without matching arbitrary family labels across candidates, we report
\begin{equation}
 P_{ij}=\sum_{m=1}^{M}W_m S_{ij}(Z_m),
 \qquad Z_\star=Z_{\arg\max_m W_m}.
 \label{eq:pairprob}
\end{equation}
Here $P_{ij}$ is the pairwise same-family probability under the candidate weights, and $Z_\star$ is the highest-weight assignment.
Each candidate is a consistent partition; thresholding $P_{ij}$ independently need not produce one.
Also, $\sum_jP_{ij}$ is an expected family size, not a quantity constrained to equal one.

\section{Computational Inference}
\label{sec:computation}
Enumerating all partitions becomes impractical as the catalog grows, so we use inexpensive affinities to propose candidates before evaluating Equation~\ref{eq:score}.
Catalog-based logistic affinities and optional source-plane affinities generate partitions through thresholding, clustering, and seeded family growth.
All-singleton and all-in-one assignments are included.
Family labels are canonicalized and duplicate partitions removed, so repeated proposals do not receive extra probability.
Appendix~\ref{app:proposals} gives the proposal equations.

After scoring, the pair probabilities can serve as soft targets for another logistic fit and candidate search.
This is iterative proposal refinement, not sampling from a proven posterior transition rule.
An optional repair pass expands the highest-scoring candidates through a family merge, a single-image reassignment, or splitting an image into a singleton.
The expanded set is deduplicated, rescored, and renormalized.
Back-projected positions are cached across families, and repeated family scores are reused.
The main calculation holds the lens model fixed; Appendix~\ref{app:optional} describes a separate implementation that averages over supplied lens realizations.

\section{Validation}
\label{sec:validation}
Assignment recovery and probability calibration require different tests.
For small catalogs, exhaustive enumeration can establish whether the candidate search finds the largest scores and reproduces the pair probabilities obtained from all partitions under the same scoring rule.
Mock catalogs with hidden family labels can then measure exact recovery, missed counterparts, and incorrect merges.
\todo{Insert catalog sizes, mock construction, recovery measurements, and a figure comparing true and inferred assignments.}

Calibration requires repeated mocks that compare reported $P_{ij}$ with the fraction of genuinely associated pairs in each probability bin.
Uncertainty estimates must account for pairs sharing images or families.
Tests should vary lens-model errors, redshift uncertainty, missing bands, and unrelated objects, then repeat the analysis with different candidate-search budgets, temperatures, and structural penalties.
\todo{Insert calibration curves, candidate-completeness measurements, sensitivity results, and runtime measurements.}

\section{Discussion and Conclusions}
\label{sec:discussion}
Weighting complete assignments retains competing family interpretations while preventing an image from occupying two families in any one candidate.
The derived pair probabilities make this ambiguity visible without requiring family labels to agree across candidates.
The method provides a way to rank alternative image catalogs for subsequent lens modeling.

Its current scores combine marginalized catalog components with a profiled lensing term and empirical preferences.
An inaccurate lens model can favor the wrong association, and a restricted candidate set can hide alternatives.
Correlated colors, lens-dependent morphology, and repeated use of catalog evidence also limit a posterior interpretation.
Replacing these approximations with a joint generative model would be a further development.
For the implemented method, the next scientific step is to establish recovery and calibration under controlled departures from its assumptions.

\section*{Data and Code Availability}
\todo{Specify the public software release, input catalogs, validation products, and persistent links.}

\begin{acknowledgments}
I.D. acknowledges support from NASA under award No. 80NSSC25K0311 under the NASA FINESST program.
P.N. gratefully acknowledges funding from the Department of Energy grant DE-SC0017660; support from the John Templeton Foundation via grant 126613 and from the NASA GLIMPSE JWST GO-03293.026 for this work.
\end{acknowledgments}

\appendix
\section{Gaussian Catalog Integral}
\label{app:gaussian}
The shared-value integral can be evaluated without sampling a separate latent parameter for each family.
Suppress the component index $a$, write $n=|F_a|$, and define $d_i=y_i-m$ and
\begin{equation}
 \begin{aligned}
 Q&=\tau^{-2}+\sum_i v_i^{-1},\\
 h&=\sum_i d_i/v_i,\qquad J=\sum_i d_i^2/v_i.
 \end{aligned}
 \label{eq:gaussstats}
\end{equation}
Here $Q$ is the total precision, $h$ the weighted offset, and $J$ the weighted squared offset; sums run over valid measurements.
For $\delta=\mu-m$, the exponent in Equation~\ref{eq:catalogmarginal} follows from
\begin{equation}
 \frac{\delta^2}{\tau^2}+\sum_i\frac{(d_i-\delta)^2}{v_i}
 =Q(\delta-h/Q)^2+J-h^2/Q.
 \label{eq:complete}
\end{equation}
Integrating the square yields the normalized density
\begin{equation}
 E=(2\pi)^{-n/2}
       \left(\tau^2Q\prod_i v_i\right)^{-1/2}
       \exp\!\left[-\frac12(J-h^2/Q)\right].
 \label{eq:closed}
\end{equation}
Thus $E$ and $B$ have the same units and their ratio is dimensionless.
The conditions $v_i>0$ and $\tau^2>0$ make the integral finite.
Equations~\ref{eq:gaussstats}--\ref{eq:closed} describe the analytic calculation; the fewer-than-two convention is applied separately.

The lensing calculation instead optimizes the shared position.
Differentiating the logarithm of Equation~\ref{eq:lensjoint} gives
\begin{equation}
 \boldsymbol{\nabla}_{\boldsymbol\beta_F}\log\mathcal L_F
 =\sum_{i\in F}\frac{\widehat{\boldsymbol\beta}_i
                         -\boldsymbol\beta_F}{u_i}.
 \label{eq:centroidderivative}
\end{equation}
Its zero is the centroid in Equation~\ref{eq:profile}; its Hessian is $-(\sum_i u_i^{-1})\boldsymbol I_2$, which is negative definite for positive variances.
No integral over the source position is performed.

\section{Catalog Inputs and Background Densities}
\label{app:catalogdetails}
To remove a common magnitude offset, the code centers each image over its available band set $\mathcal B_i$,
\begin{equation}
 c_{ib}=m_{ib}-\frac{1}{|\mathcal B_i|}
                   \sum_{r\in\mathcal B_i}m_{ir},
 \label{eq:colors}
\end{equation}
where $m_{ib}$ is the measured magnitude in band $b$.
At least two bands are required.
Different band sets imply different centering offsets, and even a common set produces correlated $c_{ib}$.
The implemented variances are $s_{z,i}^2$ for redshift, a fixed $s_c^2$ for each centered magnitude, and fixed $s_h^2$ for each morphology component.
The historical morphology option uses the same numerical width for size in arcseconds and dimensionless ellipticity.
These are separate dimensional quantities, not a common physical scatter.

Population parameters are estimated from the full input catalog.
Let $\overline y_a$ and $S_a$ be the finite-sample mean and population-form standard deviation (divisor equal to the sample size) for component $a$.
For redshift, $m_z=b_z$ is the median positive redshift and $\tau_z=t_z=\max(3,S_z)$.
With no valid redshift the location is 2; a single value has $S_z=0$.
For colors and morphology, $m_a=b_a=\overline y_a$, $\tau_a=\max(5s_a,1)$, and $t_a=\max(S_a,s_a)$ in that component's catalog units.
Missing means are replaced by zero; a missing or zero $S_a$ uses $\tau_a$ before applying the width floor.
These data-dependent choices are empirical hyperparameters, not external population measurements.

The lens background describes the spread of the entire back-projected catalog rather than a proposed family.
For coordinate $d\in\{1,2\}$ at grid point $g$, let $\overline\beta_{gd}$ and $S_{gd}^2$ be its finite-catalog mean and variance.
The cached background is
\begin{equation}
 \begin{aligned}
 p_{\bg}(\boldsymbol{\beta}\mid z_g)
 &=\prod_{d=1}^{2}\N(\beta_d\mid\overline\beta_{gd},V_{gd}),\\
 V_{gd}&=\max(S_{gd}^2,\sigma_{\rm pos}^2,
                          10^{-8}\,\mathrm{arcsec}^2),
 \end{aligned}
 \label{eq:lensbackground}
\end{equation}
where $\sigma_{\rm pos}$ is the fixed default positional scale supplied when constructing the cache.
It neglects covariance between the coordinates and uses the same catalog being assigned.
Nonfinite positions are omitted; a family with fewer than two valid source positions receives a nonfinite lens score unless it is a singleton.

\section{Proposal Probabilities and Search}
\label{app:proposals}
\subsection{Catalog affinities and iterative fitting}
Fast pair comparisons guide the search without adding a pair likelihood to the final score.
Let $\boldsymbol f_{ij}$ be a dimensionless feature vector and $\boldsymbol\theta$ its logistic coefficients.
The proposal affinity is
\begin{equation}
 q_{ij}=\{1+\exp[-\boldsymbol f_{ij}^{\mathsf T}
                              \boldsymbol\theta]\}^{-1}.
 \label{eq:logistic}
\end{equation}
The feature vector contains an intercept, proximity, radial similarity, opposite-side geometry, redshift and color similarity, size and ellipticity similarity, mean reliability, and a shared-provenance indicator.
For angular separation $d_{ij}$, radii $R_i=\|\boldsymbol x_i\|$, polar angles $\phi_i$ in radians, and $R_0=\max[\operatorname{median}(R_i),1\,\mathrm{arcsec}]$, the geometric features are
\begin{equation}
 \begin{gathered}
 \frac{1}{1+d_{ij}/(10\,\mathrm{arcsec})},\qquad
 e^{-|R_i-R_j|/R_0},\\
 \frac{1-\cos(\phi_i-\phi_j)}{2}.
 \end{gathered}
 \label{eq:geofeatures}
\end{equation}
They refer to the adopted coordinate origin, not an inferred lens center.
The remaining similarities are
\begin{equation}
 \begin{aligned}
 &\exp\!\left[-\frac{(\hat z_i-\hat z_j)^2}
                         {2(s_{z,i}^2+s_{z,j}^2)}\right],\quad
   e^{-D_{c,ij}/(1\,\mathrm{mag})},\\
 &e^{-|r_i-r_j|/(1\,\mathrm{arcsec})},\qquad
   e^{-|e_i-e_j|},
 \end{aligned}
 \label{eq:similarities}
\end{equation}
where $r_i$ is image size, $e_i$ is ellipticity, and
\begin{equation}
 D_{c,ij}^2=\frac{1}{|\mathcal B_{ij}|}
       \sum_{b\in\mathcal B_{ij}}(c_{ib}-c_{jb})^2
 \label{eq:colordistance}
\end{equation}
uses the common valid band set $\mathcal B_{ij}$, with at least two bands.
Missing similarities default to 0.5, except missing separation uses $100\,\mathrm{arcsec}$.
The reliability feature is $(\rho_i+\rho_j)/2$, where $\rho_i\in[0,1]$ is the supplied image reliability (default 1); the provenance indicator is one for equal, nonempty catalog-source labels.

For targets $t_{ij}\in[0,1]$ and fitting weights $v_{ij}^{\rm fit}$, the conditional Bernoulli objective and Gaussian coefficient penalty are
\begin{equation}
 \begin{aligned}
 \mathcal L_{\rm fit}&=\prod_{i<j}
   \left[q_{ij}^{t_{ij}}(1-q_{ij})^{1-t_{ij}}\right]^{v_{ij}^{\rm fit}},\\
 \boldsymbol\theta_\star
 &=\arg\min_{\boldsymbol\theta}
   \left[-\log\mathcal L_{\rm fit}
         +\frac{\|\boldsymbol\theta\|^2}{2s_\theta^2}\right].
 \end{aligned}
 \label{eq:fit}
\end{equation}
The fixed $s_\theta$ is the coefficient prior width.
For fractional targets this is a weighted cross-entropy objective, not independent observations of all pair labels.
After a scoring iteration, unanchored pairs use $t_{ij}=P_{ij}$ and $v_{ij}^{\rm fit}=1$; anchors use $t_{ij}=a_{ij}$ and $v_{ij}^{\rm fit}=\max(r_{ij},1)$.
In proposal matrices, anchor affinities are overwritten by $1-\epsilon$ or $\epsilon$, with $\epsilon=10^{-6}$.

Known equal family labels create positive anchors of strength 8.
Redshift differences of at least 0.8 and at least three combined standard deviations create negative anchors of strength 12, overriding a positive anchor.
Otherwise-unanchored pairs with at least three common bands and $D_{c,ij}\geq1\,\mathrm{mag}$ receive negative anchors of strength 5.
These finite penalties may disagree; they do not guarantee compatibility of all anchors.

\subsection{Lens affinities and candidate construction}
Images that nearly meet in the source plane can suggest candidates missed by catalog similarities.
Using the pair redshift weights from Equation~\ref{eq:redshift}, the lens affinity is
\begin{equation}
 H_{ij}=\sum_g\pi_{\{i,j\}g}
 \exp\!\left[-\frac{\|\widehat{\boldsymbol\beta}_i(z_g)
                    -\widehat{\boldsymbol\beta}_j(z_g)\|^2}
 {2(\sigma_i^2+\sigma_j^2+\sigma_{\rm src}^2)}\right].
 \label{eq:lensaffinity}
\end{equation}
The single scatter term in this proposal denominator follows the implementation; it is not the variance of a difference of the two independent Gaussians in Equation~\ref{eq:lensjoint}.
An optional combined affinity satisfies
\begin{equation}
 \logit q_{ij}^{\rm comb}=\logit q_{ij}
                          +w_H\logit H_{ij},
 \label{eq:combinedaffinity}
\end{equation}
where $\logit p=\log[p/(1-p)]$ and $w_H$ is a fixed proposal weight.
Affinities are clipped to $[\epsilon,1-\epsilon]$ before taking logits.

Thresholds between 0.05 and 0.95 generate connected components, sequential complete-link groups, and agglomerative groups using minimum or mean cross-group affinity.
Anchor-seeded growth merges the highest-mean-affinity compatible groups.
Pivot proposals attach unassigned images when their affinity logit to the current pivot exceeds a threshold, using several image orderings.
Randomized proposals use entries
\begin{equation}
 \widetilde q_{ij}\sim\operatorname{Beta}
 \left(\max(20q_{ij},10^{-3}),
       \max[20(1-q_{ij}),10^{-3}]\right)
 \label{eq:noisy}
\end{equation}
and thresholds drawn uniformly between 0.25 and 0.75.
Here $\operatorname{Beta}$ is the beta distribution with the two stated shape parameters.
Initial candidates also undergo single-image splits, moves, and family merges.

Optional source-plane proposals link images within an angular radius at a grid redshift, retaining images whose individual grid weights exceed a chosen threshold.
For a resulting group $F$, the proposal ranking is
\begin{equation}
 h(F)=\log\sum_g\pi_{Fg}e^{U_F(z_g)}
      +w_B\sum_{\substack{i<j\\i,j\in F}}\logit q_{ij}
      +\log|F|,
 \label{eq:beamscore}
\end{equation}
where $U_F$ is the unnormalized source score defined in Appendix~\ref{app:optional} and $w_B$ weights catalog affinities for this search only.
A bounded search combines nonoverlapping groups by adding $h(F)$, retaining only the best state for a given set of used images; unassigned images become singletons.
All retained partitions are evaluated with Equation~\ref{eq:score} afterward.

\section{Optional Source Scoring Components}
\label{app:optional}
\subsection{Reliability mixture}
A broad residual component can reduce the influence of an unreliable image.
At the centroid from Equation~\ref{eq:profile}, let $d_i^2=\|\widehat{\boldsymbol\beta}_i-\widehat{\boldsymbol\beta}_F\|^2$ and let $u_{\rm out}=\sigma_{\rm out}^2$ be a fixed broad variance.
The optional product is
\begin{equation}
 \begin{aligned}
 \mathcal L_F^{\rm mix}=\prod_{i\in F}\Big[&
 \rho_i\N_2(\widehat{\boldsymbol\beta}_i\mid
               \widehat{\boldsymbol\beta}_F,u_i\boldsymbol I_2)\\
 {}+&(1-\rho_i)\N_2(\widehat{\boldsymbol\beta}_i\mid
               \widehat{\boldsymbol\beta}_F,u_{\rm out}\boldsymbol I_2)
 \Big].
 \end{aligned}
 \label{eq:mixture}
\end{equation}
Both components share the fitted centroid; the code does not reoptimize it for the mixture.
The cached likelihood-ratio path retains the Gaussian prefactors, whereas the raw path evaluates
\begin{equation}
 U_F(z)=\sum_{i\in F}\log\left[
 \rho_i e^{-d_i^2/(2u_i)}
 +(1-\rho_i)e^{-d_i^2/(2u_{\rm out})}\right].
 \label{eq:raw}
\end{equation}
This raw score is not a normalized density.
Numerically, each mixture coefficient is clipped separately to $[10^{-6},1]$ without renormalization, so Equation~\ref{eq:mixture} describes its interior-weight form.
The main-text Gaussian derivation corresponds to omitting this optional mixture.

\subsection{Averaging over supplied lens realizations}
Lens uncertainty can change family rankings, so a separate callback accepts multiple ray-shooting models.
For each family it uses the inverse-variance mean redshift $\bar z_F=(\sum_{i\in F_z}s_{z,i}^{-2})^{-1}\sum_{i\in F_z}\hat z_i s_{z,i}^{-2}$.
This callback evaluates the raw score rather than the cached redshift-grid likelihood ratio.
For realization $r$, let $u_r(Z)=\sum_{F\in Z}U_F(\bar z_F;\boldsymbol\eta_r)$ and let $\mathcal R_Z$ contain its finite values.
The implemented average is
\begin{equation}
 U_{\rm avg}(Z)=\log\left[
 \frac{1}{|\mathcal R_Z|}\sum_{r\in\mathcal R_Z}e^{u_r(Z)}\right].
 \label{eq:lensaverage}
\end{equation}
An empty $\mathcal R_Z$ gives $-\infty$.
This averages scores over the supplied models, with a partition-dependent denominator if some fail.
It does not automatically marginalize the main-text composite score or establish joint inference of partitions and lens parameters.


\bibliography{refs}{}
\bibliographystyle{aasjournalv7}

\end{document}

